% The national range of the name Mandji: ten blocks laid end to end from 1000 (Eq. 6), each with its
% numbers issued at first issue, the rest of its main part and its tail reserve (Eq. 9), then the spare.
% Data: mandji-blocks.dat, written by tools/paper-figures.py.
\documentclass[tikz,border=3pt]{standalone}
\usepackage{fontspec}
\setmainfont{LibertinusSans}[Extension=.otf,UprightFont=*-Regular,BoldFont=*-Bold,ItalicFont=*-Italic]
\usepackage{unicode-math}
\setmathfont{LibertinusMath-Regular.otf}
\usepackage{pgfplots}
\pgfplotsset{compat=1.18}
% English by default; the French paper compiles <name>-fr.tex, which sets \fr to 1 and inputs this file.
\providecommand{\fr}{0}
\newcommand{\T}[2]{\ifnum\fr=1 #2\else #1\fi}
\ifnum\fr=1 \pgfkeys{/pgf/number format/.cd, use comma, 1000 sep={\,}}\fi
\definecolor{ink}{HTML}{1F5FA8}\definecolor{flaggreen}{HTML}{009E60}\definecolor{yellow}{HTML}{FCD116}
\definecolor{flagblue}{HTML}{3A75C4}\definecolor{amber}{HTML}{D9A400}
\usepackage{pgfplotstable}
\usetikzlibrary{patterns,decorations.pathreplacing}
\begin{document}
\begin{tikzpicture}[x=1cm,y=1cm]
\pgfplotstableread{mandji-blocks.dat}\blocks
\pgfplotstablegetrowsof{\blocks}\let\nrows\pgfplotsretval
% whole range 1000-9999 on 14 cm
\def\W{14}
\pgfmathsetmacro{\sx}{\W/9000}
\fill[gray!12] (0,2.8) rectangle (\W,3.4);
\fill[flagblue!60] (0,2.8) rectangle ({(4210-1000)*\sx},3.4);
\draw[gray] (0,2.8) rectangle (\W,3.4);
\draw[flaggreen, thick] ({(9100-1000)*\sx},2.7) -- ({(9100-1000)*\sx},3.5);
\node[font=\footnotesize, anchor=south] at ({(4210-1000)*\sx/2},3.4) {\T{ten blocks, 1000–4209}{dix blocs, 1000–4209}};
\node[font=\footnotesize, anchor=south] at ({(6650-1000)*\sx},3.4) {\T{spare for new and extension blocks}{réserve pour les nouveaux blocs et les extensions}};
\node[font=\footnotesize, anchor=south, align=center, text=flaggreen] at ({(9550-1000)*\sx},3.4) {\T{top tenth}{dixième supérieur}};
\foreach \t in {1000,2000,...,9000} {\draw[gray] ({(\t-1000)*\sx},2.8) -- ({(\t-1000)*\sx},2.7) node[below, font=\scriptsize] {\t};}
\draw[gray] (\W,2.8) -- (\W,2.7) node[below, font=\scriptsize] {9999};
% zoom on 1000-4209
\pgfmathsetmacro{\zx}{\W/3210}
\draw[gray!60, dashed] (0,2.25) -- (0,1.0);
\draw[gray!60, dashed] ({(4210-1000)*\sx},2.25) -- (\W,1.0);
\foreach \i in {0,...,9} {
  \pgfplotstablegetelem{\i}{lo}\of\blocks \let\lo\pgfplotsretval
  \pgfplotstablegetelem{\i}{hi}\of\blocks \let\hi\pgfplotsretval
  \pgfplotstablegetelem{\i}{mainhi}\of\blocks \let\mh\pgfplotsretval
  \pgfplotstablegetelem{\i}{usedhi}\of\blocks \let\uh\pgfplotsretval
  \pgfplotstablegetelem{\i}{area}\of\blocks \let\ar\pgfplotsretval
  \pgfplotstablegetelem{\i}{n}\of\blocks \let\nn\pgfplotsretval
  \pgfmathsetmacro{\a}{(\lo-1000)*\zx}
  \pgfmathsetmacro{\u}{(\uh+1-1000)*\zx}
  \pgfmathsetmacro{\m}{(\mh+1-1000)*\zx}
  \pgfmathsetmacro{\b}{(\hi+1-1000)*\zx}
  \fill[flagblue] (\a,0.4) rectangle (\u,1.0);
  \fill[flagblue!30] (\u,0.4) rectangle (\m,1.0);
  \fill[pattern=north east lines, pattern color=amber] (\m,0.4) rectangle (\b,1.0);
  \draw[white, line width=0.8pt] (\a,0.4) -- (\a,1.0);
  \pgfmathsetmacro{\mid}{(\a+\b)/2}
  \pgfmathtruncatemacro{\odd}{mod(\i,2)}
  \ifnum\odd=0 \node[font=\tiny, anchor=north, align=center] at (\mid,0.38) {\ar\\\lo};
  \else \node[font=\tiny, anchor=south, align=center] at (\mid,1.02) {\ar\\\lo}; \fi
}
\draw[gray] (0,0.4) rectangle (\W,1.0);
% legend
\fill[flagblue] (0,-0.75) rectangle (0.35,-0.5); \node[font=\footnotesize, anchor=west] at (0.4,-0.62) {\T{numbers issued}{numéros émis}};
\fill[flagblue!30] (3.4,-0.75) rectangle (3.75,-0.5); \node[font=\footnotesize, anchor=west] at (3.8,-0.62) {\T{rest of the main part (infill)}{reste de la partie principale (insertion)}};
\fill[pattern=north east lines, pattern color=amber] (8.6,-0.75) rectangle (8.95,-0.5); \draw[gray] (8.6,-0.75) rectangle (8.95,-0.5);
\node[font=\footnotesize, anchor=west] at (9.0,-0.62) {\T{tail reserve (last third)}{réserve de queue (dernier tiers)}};
\end{tikzpicture}
\end{document}
